% Part 1 – Preface, Course Information and Introduction

% ---------------------------------------------------------------------------
% This file constitutes the first logical section of the full report.  It
% introduces the assignment context, provides background on the motivation for
% studying generative models, and lays out the overall structure of the report.
% The content is designed to be self–contained and approachable, making it a
% natural starting point for readers who may not be familiar with energy–
% based or score–based generative modeling.  Throughout the text we adopt a
% narrative tone, guiding the reader through the rationale for the project
% while retaining scientific rigor.  Colour and emphasis commands defined in
% the main document are used consistently to highlight key concepts without
% overwhelming the reader.

\section{Preface and Motivation}

\subsection{University and Course Context}

The work presented here was carried out as part of the \emph{Trustworthy~AI}
course taught at the \textit{School of Electrical and Computer Engineering},
\textbf{University of Tehran}.  The assignment forms the third homework in
the course and is intended to deepen students' understanding of modern
generative modeling techniques, including \emph{Energy–Based Models}
(EBMs) and \emph{Noise Conditional Score Networks} (NCSNs).  The problem
specification, provided in Persian and translated in the project description,
outlines both theoretical questions and an extensive implementation task
involving the MNIST dataset.  Students are expected to deliver both a working
Python codebase and a comprehensive report documenting their methodology and
findings.  This document constitutes the latter.

\subsection{Invocation and Dedication}

In keeping with local tradition and the formal tone of academic reports in
Iran, we begin by invoking a short dedication.  The original assignment
document opens with the phrase ``\emph{In the Name of God}'' to place the
scientific endeavour within a broader cultural and ethical context.  We
include the dedication here to honour that tradition while emphasising that
the core of this work remains firmly grounded in empirical investigation and
mathematical reasoning.

\subsection{Assignment Overview}

The assignment consists of two broad components: a set of theoretical
questions exploring the properties of EBMs and score–based generative models,
and a practical component requiring the implementation of these models on the
MNIST handwritten digits dataset.  The theoretical questions probe topics
such as the role of the partition function in energy–based modeling, the
challenges of sampling and normalisation in high dimensions, and the
fundamental principles of denoising score matching and multi–scale noise
perturbation.  The practical component requires constructing a convolutional
energy network and training it via contrastive divergence, implementing
Langevin sampling for generation and denoising, and building a U–Net based
score network that can generate digits via annealed Langevin dynamics.  A
conditional version of the score network must also be trained to allow class
controlled generation.  In both cases students are asked to report training
curves, display representative samples, perform denoising experiments at
multiple noise levels, and analyse the factors influencing sample quality.

\subsection{Why Generative Modeling?}

Generative modeling has emerged as a cornerstone of contemporary machine
learning.  Unlike discriminative models, which learn to map inputs to labels
or real–valued predictions, generative models aim to learn the entire
distribution of the data.  Once trained, a generative model can synthesise
novel samples that resemble the training examples, interpolate between
existing examples, perform reconstruction and denoising, or serve as a
regulariser for downstream tasks.  Classical approaches to generative
modeling include \emph{Gaussian Mixture Models}, \emph{autoregressive
models}, \emph{Variational Autoencoders}, and \emph{Generative Adversarial
Networks}.  More recent methods, such as \emph{diffusion models} and
\emph{score–based generative models}, have demonstrated remarkable
performance on high–dimensional data like natural images.

Within this landscape, EBMs and NCSNs occupy an interesting niche.  EBMs
define an unnormalised density via an energy function and avoid the need to
compute an explicit partition function.  Score–based models, on the other
hand, model the gradient of the log–density and thereby sidestep
normalisation entirely.  Both approaches harness stochastic differential
equations—most notably Langevin dynamics—to transform random noise into
coherent samples.  Understanding these models therefore provides insight
into a broad class of generative techniques that emphasise sampling and
implicit densities over explicit likelihood evaluation.

\subsection{Report Structure}

To make the presentation clear and modular, this report has been divided
into eight parts, each compiled as a separate \TeX{} file and included into
the final document.  Each part is designed to be self–contained and
approximately two hundred lines in length, providing detailed explanations,
mathematical derivations, implementation details, and experimental analyses.
The division into multiple files allows us to maintain a logical flow
without overwhelming the reader with a monolithic block of text.  It also
facilitates the reuse of macros and colour definitions introduced in the
main preamble.

\begin{enumerate}[leftmargin=*,label=\arabic*.]
  \item \textbf{Preface and Motivation (this section).}  This part lays out
    the motivation for the project, summarises the assignment instructions
    provided in the original prompt, and outlines the structure of the
    report.  It introduces generative modeling concepts at a high level and
    explains why EBMs and NCSNs are important to study.  It also places the
    assignment within the context of the course on \emph{Trustworthy AI} at
    the University of Tehran.
  \item \textbf{Energy–Based Models: Theory.}  The second part delves into
    the theoretical foundations of EBMs, including the definition of the
    energy function, the role of the partition function, derivations of the
    log–likelihood gradient, and sampling methods such as Langevin dynamics
    and rejection sampling.  The mixture–of–Gaussians example is revisited
    to illustrate why simple gradient–based sampling may fail to capture
    multimodal distributions.  Key challenges such as computing the
    partition function and drawing samples in high dimensions are analysed.
  \item \textbf{Energy–Based Models: Implementation.}  Having covered the
    theory, the third part describes how to implement an EBM for MNIST.
    Details include data preprocessing, the convolutional architecture
    specified in Table~\ref{tab:ebm-arch}, the training loop using contrastive
    divergence, Langevin sampling for both positive and negative phases, and
    practical considerations such as regularisation and hyperparameter
    selection.  Pseudocode is provided and key functions are listed to
    facilitate reproducibility.
  \item \textbf{Energy–Based Models: Results and Analysis.}  This part
    presents training curves, samples generated at different stages of
    training, and denoising experiments at multiple noise levels.  Each
    figure is accompanied by detailed commentary explaining the observed
    behaviour and relating it to the underlying theory.  We analyse the
    influence of the number of Langevin steps, step size, regularisation
    strength, and other hyperparameters on sample quality.  The limitations
    of the model and suggestions for improvement are discussed.
  \item \textbf{Score–Based Models: Theory.}  The fifth part introduces
    score–based generative modeling.  We define the score function,
    demonstrate its independence from the partition function, and derive the
    denoising score matching objective.  The difficulties of computing
    divergence terms in the original score matching formulation are
    explained, and multi–scale noise perturbation is presented as a remedy.
    We also revisit the mixture–of–Gaussians example to show why the score
    ignores mixing weights.  Connections to stochastic differential
    equations and diffusion processes are highlighted.
  \item \textbf{Score–Based Models: Implementation.}  This section outlines
    the practical implementation of Noise Conditional Score Networks.  The
    architecture, inspired by U–Nets and equipped with FiLM conditioning
    based on Gaussian Fourier embeddings, is described in detail.  We cover
    the weighted denoising score matching loss, the geometric noise ladder,
    training schedules, and annealed Langevin dynamics for sampling.  A
    class–conditional extension is developed by adding a label embedding and
    training the network on labelled MNIST data.
  \item \textbf{Score–Based Models: Results and Analysis.}  As with the
    EBM, we report training curves, show samples produced by the score
    network, visualise the annealing trajectory from noise to digits, and
    evaluate denoising across several noise levels.  The conditional model
    is used to generate a grid of digits conditioned on each class, and we
    assess fidelity and diversity.  Comparisons with the EBM results are
    drawn.
  \item \textbf{Discussion, Limitations and Conclusions.}  The final part
    brings together the theoretical insights and experimental findings from
    the preceding sections.  It contrasts the strengths and weaknesses of
    EBMs and NCSNs, summarises the hyperparameter sensitivity analyses,
    enumerates limitations of the current implementations, and suggests
    directions for future work.  Ethical considerations are also briefly
    discussed.  The report concludes with key takeaways and reflections on
    the learning journey.
\end{enumerate}

\subsection{Reading Guide}

The report is deliberately long and detailed because it aims to serve as a
self–contained reference for students and researchers interested in
understanding energy–based and score–based generative models at both
conceptual and practical levels.  To facilitate navigation, each part begins
with its own introduction and ends with a summary of key points.  Readers
who are already familiar with the underlying theory may wish to skim
Part~2 and Part~5 and focus on the implementation and results sections.

\subsection{Notation and Conventions}

Throughout this report, vectors are denoted in boldface (e.g., $\mathbf{x}$)
and matrices in uppercase boldface (e.g., $\mathbf{W}$).  The symbol $\theta$
collectively denotes the parameters of a model.  Energy functions are
written as $E_\theta(\mathbf{x})$ or $f_\theta(\mathbf{x})$ interchangeably, and
probability densities defined via energy functions are written as
$p_\theta(\mathbf{x}) \propto \exp\big(-f_\theta(\mathbf{x})\big)$.  The gradient
with respect to a vector argument is denoted by $\nabla$, and the
divergence of a vector field $\mathbf{s}(\mathbf{x})$ is written as
$\nabla \cdot \mathbf{s}(\mathbf{x}) = \sum_i \partial s_i(\mathbf{x}) / \partial x_i$.
For clarity, we consistently refer to the training distribution as
$p_{\text{data}}$ and to the model distribution as $p_\theta$.

In citations, we use bracket notation to refer to specific ranges of lines in
source documents.  For example, the statement that EBMs do not require
proper normalisation, thus providing more flexibility in designing loss
functions, is supported by LeCun et~al. who note that probabilistic models
must be normalised but ``\emph{since EBMs have no requirement for proper
normalisation, this problem is naturally circumvented}''【861141214805006†L26-L33】.  This
property underlies many of the advantages of energy–based modeling and is
revisited throughout the report.

\subsection{A Word on Scientific Rigor}

The instructions emphasise that all explanations should be scientific and
supported by reliable sources.  Where possible, we cite primary literature
and authoritative tutorials, such as the classic work by LeCun et~al. on
energy–based learning【861141214805006†L26-L33】.  For diffusion and score–based models,
we rely on established theoretical results and the seminal papers by Song
and collaborators.  While network constraints may limit access to some
online resources, the core mathematical derivations are self–contained and
draw upon widely known results in statistical physics and stochastic
processes.  We strive to maintain transparency and reproducibility in both
our theoretical arguments and implementation details.

\subsection{Conclusion of Part 1}

This introductory part has set the stage for the remainder of the report.  We
have situated the assignment within the educational context of a graduate
course on \emph{Trustworthy AI}, summarised the tasks to be accomplished,
explained why generative modeling is worth studying, and described how the
report is organised.  In the following part we will delve into the theory
of energy–based models, starting from first principles and building up to
the maximum–likelihood gradient used to train the MNIST EBM.  Readers who
are eager to see code and results may jump ahead to the implementation
sections, but we encourage a thorough reading of the theory to fully
appreciate the challenges and solutions presented later on.